\documentclass{article}
\usepackage[utf8]{inputenc}
\usepackage{amsmath}
\usepackage{graphicx}

\title{{{TITLE}}}
\author{{{AUTHOR}}}
\date{\today}

\begin{document}

\maketitle

\section{Introduction}

This is an example project. Edit the source on the left; the PDF on the right
recompiles automatically. You can cite references from \texttt{references.bib},
such as \TeX's reference manual~\cite{knuth1984}, and refer to numbered things
such as Equation~\ref{eq:energy}, Table~\ref{tab:results}, and
Figure~\ref{fig:example}.

\section{Mathematics}

Inline mathematics looks like $e^{i\pi} + 1 = 0$, and aligned equations look
like this:
\begin{align}
  E &= mc^2 \label{eq:energy} \\
  \int_0^1 x^2 \, dx &= \frac{1}{3}
\end{align}

\section{Tables and figures}

\begin{table}[h]
  \centering
  \begin{tabular}{lrr}
    \hline
    Method & Accuracy & Time (s) \\
    \hline
    Baseline & 0.81 & 12 \\
    Improved & 0.93 & 15 \\
    \hline
  \end{tabular}
  \caption{Some example results.}
  \label{tab:results}
\end{table}

\begin{figure}[h]
  \centering
  \includegraphics[width=0.5\linewidth]{example-image}
  \caption{An example image from the \texttt{mwe} package.}
  \label{fig:example}
\end{figure}

\section{Conclusion}

\LaTeX{} itself is described by Lamport~\cite{lamport1994}.

\bibliographystyle{plain}
\bibliography{references}

\end{document}
